%! Author = sriadityavedantam
%! Date = 3/30/26

\section{Chapter 7}

\begin{problem}{
    Let $R$ be a ring with identity and let $I$ be an ideal in $R$.

    \textbf{(a).}
    If $1_R \in I$, prove that $I = R$.

    \textbf{(b).}
    If $I$ contains a unit, prove that $I = R$.
}
    \textbf{(a).}
    Suppose $1_R \in I$.
    Let $r\in R$ be arbitrary.
    Since $I$ is an ideal and $1_R\in I$, we have
    \[
        r1_R = r \in I.
    \]
    Thus every element of $R$ lies in $I$.
    Hence $I = R$.

    \textbf{(b).}
    Suppose $I$ contains a unit $u\in R$.
    Then there exists $u^{-1}\in R$ such that
    \[
        u^{-1}u = 1_R.
    \]
    Since $u\in I$ and $I$ is an ideal, we have
    \[
        u^{-1}u = 1_R \in I.
    \]
    By part \textbf{(a)}, it follows that $I = R$.
\end{problem}
\begin{problem}{
    If $I$ is an ideal in a field $\mathbb{F}$, prove that
    \[
        I = (0)_{\mathbb{F}}
        \quad\text{or}\quad
        I = \mathbb{F}.
    \]
}
    If $I = (0)_{\mathbb{F}}$, then we are done.
    Suppose instead that $I\neq (0)_{\mathbb{F}}$.
    Then there exists some $a\in I$ with $a\neq 0$.
    Since $\mathbb{F}$ is a field, $a$ is a unit.
    Therefore, by the previous problem, $I = \mathbb{F}$.
    Hence every ideal in a field is either $(0)_{\mathbb{F}}$ or $\mathbb{F}$.
\end{problem}
\begin{problem}{
    Let $J$ be an ideal in $R$.
    Prove that $I$ is an ideal, where
    \[
        I \coloneqq \{r\in R : rt = 0_R \text{ for all } t\in J\}.
    \]
}
    First, note that $0_R\in I$, since for every $t\in J$,
    \[
        0_R t = 0_R.
    \]
    Now let $a,b\in I$.
    Then for every $t\in J$,
    \[
        (a-b)t = at - bt = 0_R - 0_R = 0_R.
    \]
    Thus $a-b\in I$.
    Now let $r\in R$ and $a\in I$.
    Then for every $t\in J$,
    \[
        (ra)t = r(at) = r0_R = 0_R.
    \]
    So $ra\in I$.
    Also, since $J$ is an ideal, we have $rt\in J$ for every $t\in J$.
    Therefore
    \[
        (ar)t = a(rt) = 0_R
    \]
    for every $t\in J$.
    So $ar\in I$.
    Hence $I$ is an ideal of $R$.
\end{problem}
\begin{problem}{
    Let $I$ and $J$ be ideals in $R$.
    Let $IJ$ denote the set of all possible finite sums of elements of the form $ab$, with $a\in I$ and $b\in J$, that is,
    \[
        IJ \coloneqq \{a_{1}b_1 + \hdots + a_{n}b_n : n\geq 1,\ a_k\in I,\ b_k\in J\}.
    \]
    Prove that $IJ$ is an ideal.
}
    Let
    \[
        x = \sum_{i=1}^m a_i b_i \in IJ
        \quad\text{and}\quad
        y = \sum_{j=1}^n c_j d_j \in IJ.
    \]
    Then
    \[
        x-y
        =
        \sum_{i=1}^m a_i b_i - \sum_{j=1}^n c_j d_j
        =
        \sum_{i=1}^m a_i b_i + \sum_{j=1}^n (-c_j)d_j.
    \]
    Since $I$ is an ideal, $-c_j\in I$ for each $j$.
    Thus $x-y$ is again a finite sum of terms of the form $ab$ with $a\in I$ and $b\in J$.
    So $x-y\in IJ$.
    Now let $r\in R$.
    Then
    \[
        rx = r\left(\sum_{i=1}^m a_i b_i\right) = \sum_{i=1}^m (ra_i)b_i.
    \]
    Since $I$ is an ideal, $ra_i\in I$.
    Thus $rx\in IJ$.
    Similarly,
    \[
        xr = \left(\sum_{i=1}^m a_i b_i\right)r = \sum_{i=1}^m a_i(b_i r).
    \]
    Since $J$ is an ideal, $b_i r\in J$.
    Thus $xr\in IJ$.
    Also $0\in IJ$, since if $a\in I$ and $b\in J$, then
    \[
        ab-ab = 0
    \]
    is in $IJ$.
    Therefore $IJ$ is an ideal of $R$.
\end{problem}
\begin{problem}{
    Prove that every ideal in $\mathbb{Z}$ is principal.
}
    Let $I$ be an ideal of $\mathbb{Z}$.
    If $I=\{0\}$, then
    \[
        I = (0),
    \]
    so $I$ is principal.
    Suppose $I\neq \{0\}$.
    Since $I$ contains a nonzero integer, it contains a positive integer.
    Let $c\in I$ be the smallest positive integer in $I$.
    We claim that
    \[
        I = (c).
    \]
    First, since $c\in I$ and $I$ is an ideal, every multiple of $c$ lies in $I$.
    Thus
    \[
        (c)\subseteq I.
    \]
    Now let $a\in I$.
    By the division algorithm, there exist integers $q$ and $r$ such that
    \[
        a = qc + r
        \quad\text{with}\quad
        0\leq r < c.
    \]
    Since $c\in I$, we have $qc\in I$.
    Also $a\in I$.
    Thus
    \[
        r = a-qc \in I.
    \]
    But $c$ was chosen to be the smallest positive integer in $I$.
    Since $0\leq r<c$, it follows that $r=0$.
    Hence
    \[
        a = qc \in (c).
    \]
    So $I\subseteq (c)$.
    Therefore
    \[
        I = (c),
    \]
    and every ideal in $\mathbb{Z}$ is principal.
\end{problem}
\begin{problem}{
    Let $\mathbb{F}$ be a field.
    Prove that every ideal in $\mathbb{F}[x]$ is principal.
}
    Let $I$ be an ideal in $\mathbb{F}[x]$.
    If $I=\{0\}$, then
    \[
        I = (0),
    \]
    so $I$ is principal.
    Suppose $I\neq \{0\}$.
    Choose a nonzero polynomial $b(x)\in I$ of smallest degree.
    We claim that
    \[
        I = (b(x)).
    \]
    First, since $b(x)\in I$ and $I$ is an ideal, every multiple of $b(x)$ lies in $I$.
    Thus
    \[
        (b(x))\subseteq I.
    \]
    Now let $a(x)\in I$.
    By the division algorithm in $\mathbb{F}[x]$, there exist polynomials $q(x),r(x)\in\mathbb{F}[x]$ such that
    \[
        a(x) = q(x)b(x) + r(x),
    \]
    where either $r(x)=0$ or
    \[
        \deg r(x) < \deg b(x).
    \]
    Since $b(x)\in I$, we have $q(x)b(x)\in I$.
    Also $a(x)\in I$.
    Therefore
    \[
        r(x) = a(x)-q(x)b(x)\in I.
    \]
    By the minimality of the degree of $b(x)$, we must have $r(x)=0$.
    Hence
    \[
        a(x) = q(x)b(x)\in (b(x)).
    \]
    So $I\subseteq (b(x))$.
    Therefore
    \[
        I = (b(x)),
    \]
    and every ideal in $\mathbb{F}[x]$ is principal.
\end{problem}
\begin{problem}{
    Show that every homomorphic image of a field $\mathbb{F}$ is isomorphic either to $\mathbb{F}$ itself or to the zero ring.
}
    Let $f:\mathbb{F}\to S$ be a ring homomorphism.
    Let
    \[
        K = \ker f.
    \]
    Since $\mathbb{F}$ is a field, its only ideals are $(0)$ and $\mathbb{F}$.
    Therefore
    \[
        K = (0)
        \quad\text{or}\quad
        K = \mathbb{F}.
    \]
    If $K=(0)$, then by the First Isomorphism Theorem,
    \[
        \mathbb{F}/(0)\cong \ima(f),
    \]
    so
    \[
        \ima(f)\cong \mathbb{F}.
    \]
    If $K=\mathbb{F}$, then $f$ is the zero homomorphism, so its image is the zero ring.
    Hence every homomorphic image of a field is isomorphic either to the field itself or to the zero ring.
\end{problem}
\begin{problem}{
    Use the First Isomorphism Theorem to show that
    \[
        \mathbb{Z}_{20}/(5) \cong \mathbb{Z}_5.
    \]
}
    Define
    \[
        f:\mathbb{Z}_{20}\to \mathbb{Z}_5
    \]
    by
    \[
        f([a]_{20}) = [a]_5.
    \]
    First we check that $f$ is well-defined.
    Suppose
    \[
        [a]_{20} = [b]_{20}.
    \]
    Then
    \[
        a\equiv b\mod 20,
    \]
    so
    \[
        20\mid (a-b).
    \]
    Hence
    \[
        5\mid (a-b),
    \]
    which means
    \[
        [a]_5 = [b]_5.
    \]
    Therefore $f$ is well-defined.
    Now we check that $f$ is a ring homomorphism.
    For addition,
    \begin{align*}
        f([a]_{20} + [b]_{20})
        &=
        f([a+b]_{20})\\
        &=
        [a+b]_5\\
        &=
        [a]_5 + [b]_5\\
        &=
        f([a]_{20}) + f([b]_{20}).
    \end{align*}
    For multiplication,
    \begin{align*}
        f([a]_{20}[b]_{20})
        &=
        f([ab]_{20})\\
        &=
        [ab]_5\\
        &=
        [a]_5[b]_5\\
        &=
        f([a]_{20})f([b]_{20}).
    \end{align*}
    The map is surjective, since every class $[a]_5\in \mathbb{Z}_5$ is the image of $[a]_{20}$.
    Now compute the kernel.
    We have
    \[
        f([a]_{20}) = [0]_5
    \]
    if and only if
    \[
        [a]_5 = [0]_5,
    \]
    which holds if and only if $5\mid a$.
    Thus
    \[
        \ker f = (5).
    \]
    By the First Isomorphism Theorem,
    \[
        \mathbb{Z}_{20}/(5)\cong \ima(f)=\mathbb{Z}_5.
    \]
\end{problem}
\begin{problem}{
    Let $I$ and $J$ be ideals in a ring $R$.
    Then $I\cap J$ is an ideal in $I$, and $J$ is an ideal in $I+J$.
    Prove that
    \[
        \frac{I}{I\cap J} \cong \frac{I+J}{J}.
    \]
}
    Define
    \[
        f:I\to \frac{I+J}{J}
    \]
    by
    \[
        f(a) = a+J.
    \]
    First we check that $f$ is a ring homomorphism.
    If $a,b\in I$, then
    \begin{align*}
        f(a+b)
        &=
        (a+b)+J\\
        &=
        (a+J)+(b+J)\\
        &=
        f(a)+f(b).
    \end{align*}
    Also,
    \begin{align*}
        f(ab)
        &=
        ab+J\\
        &=
        (a+J)(b+J)\\
        &=
        f(a)f(b).
    \end{align*}
    Now we show $f$ is surjective.
    An arbitrary element of $(I+J)/J$ has the form
    \[
        (a+j)+J
    \]
    with $a\in I$ and $j\in J$.
    But
    \[
        (a+j)+J = a+J = f(a).
    \]
    So $f$ is surjective.
    Now compute the kernel.
    We have
    \[
        f(a)=J
    \]
    if and only if
    \[
        a+J = J,
    \]
    which happens if and only if $a\in J$.
    Since also $a\in I$, this means
    \[
        a\in I\cap J.
    \]
    Thus
    \[
        \ker f = I\cap J.
    \]
    By the First Isomorphism Theorem,
    \[
        \frac{I}{I\cap J} \cong \frac{I+J}{J}.
    \]
\end{problem}